\section{Introduction}
Le chapitre précédent a présenté la mise en place des plots et des cultures, le pilotage de l'irrigation selon la méthode FAO-56 Penman-Monteith \cite{allen_crop_1998-1}, les outils d'acquisition ainsi que la gestion hebdomadaire et le suivi des cultures. 

Cette expérimentation a permis d'établir une base de données issue du terrain, destinée à servir de fondement à la conception d’une application capable d’automatiser ces calculs et d’appuyer le pilotage de l’irrigation. 
Ce chapitre marque donc une transition de l'expérimentation agronomique vers une conception informatique du système.

La conception du système nécessite d’abord de déterminer les éléments à automatiser, les acteurs concernés et la logique de fonctionnement. 
Pour structurer cette analyse, la méthode Merise a été retenue. Ainsi, ce chapitre présente d’abord les raisons de ce choix, puis le domaine d'étude et ses acteurs, ensuite les différents modèles de conception, notamment le MCC, le MCT et le MOT ainsi que le diagramme de Gantt. Enfin, il présente l’application issue de ces travaux.
